\section{Tính đa thức tại nhiều điểm [Cơ bản]}
\subsection{Phân tích \& Thiết kế (M0)}
    \begin{itemize}
    \item \textbf{Input:}
        \begin{itemize}
            \item File \texttt{poly.csv} chứa 1 dòng: \texttt{degree, a0, a1, ..., an} (phân tách bằng dấu phẩy).
            \item File \texttt{x\_values.csv} chứa nhiều giá trị $x$, mỗi dòng 1 giá trị.
        \end{itemize}

    \item \textbf{Xử lý:}
        \begin{itemize}
            \item \textbf{Case biên:}
            \begin{itemize}
                \item \texttt{degree = 0} $\rightarrow$ Đa thức là hằng số $P(x) = a_0$.
                \item File \texttt{x\_values.csv} rỗng $\rightarrow$ File output chỉ in ra đúng dòng header \texttt{x,P(x)}.
                \item Gặp giá trị khuyết (\texttt{NA}, \texttt{NaN}, \texttt{nan}, hoặc dòng trống) $\rightarrow$ Bỏ qua không đưa vào tính toán.
            \end{itemize}
            
            \item \textbf{Module 1 (Đọc hệ số đa thức):}
            \begin{itemize}
                \item Hàm \texttt{readPoly(path, int\&, vector<double>\&)}
                \item Tách chuỗi bằng dấu phẩy (\texttt{,}) thông qua \texttt{stringstream}.
                \item Bắt lỗi mở file (\text{fail} $\rightarrow$ \texttt{cerr} + \texttt{return false}).
            \end{itemize}

            \item \textbf{Module 2 (Đọc giá trị x):}
            \begin{itemize}
                \item Hàm \texttt{readXValues(path, vector<double>\&)}
                \item Đọc từng dòng, kiểm tra missing values.
                \item Bắt lỗi mở file.
            \end{itemize}
            
            \item \textbf{Module 3 (Tính giá trị đa thức):}
            \begin{itemize}
                \item Hàm \texttt{evalHorner(vector<double> a, double x)}
                \item Áp dụng thuật toán Horner: \texttt{res = a[n]; for i=n-1..0: res = res*x + a[i]}.
            \end{itemize}

            \item \textbf{Module 4 (Ghi output):}
            \begin{itemize}
                \item Hàm \texttt{writeOutput(path, vector<double>, vector<double>)}
                \item Ghi header \texttt{x,P(x)}. Định dạng số thực bằng \texttt{fixed} và \texttt{setprecision(6)}.
            \end{itemize}
        \end{itemize}

        \item \textbf{Output:}
        File \texttt{output.csv} với định dạng:
        \begin{lstlisting}
            x,P(x)
            <f>,<f>\end{lstlisting} 
    \end{itemize}

\subsection{Mã nguồn C++11 (Tách module)}
    \begin{lstlisting}
    #include <iostream>
    #include <fstream>
    #include <vector>
    #include <string>
    #include <sstream>
    #include <iomanip>

    using namespace std;

    // MODULE 1: Doc da thuc tu file poly.csv
    bool readPoly(string path, int& degree, vector<double>& a) {
        ifstream file(path);
        if (!file.is_open()) {
            cerr << "Khong mo duoc file " << path << "\n";
            return false;
        }

        string line;
        if (getline(file, line)) {
            stringstream ss(line);
            string token;
            
            // Doc degree
            if (getline(ss, token, ',')) {
                degree = stoi(token);
            }
            
            // Doc cac he so a0, a1, ..., an
            while (getline(ss, token, ',')) {
                // Kiem tra missing value theo quy dinh (da ngat dong de tranh tran le PDF)
                if (token == "NA" || token == "NaN" || 
                    token == "nan" || token.empty()) {
                    a.push_back(0.0); // He so khuyet xem nhu bang 0
                } else {
                    a.push_back(stod(token));
                }
            }
        }

        file.close();
        return true;
    }

    // MODULE 2: Doc cac diem x tu x_values.csv
    bool readXValues(string path, vector<double>& x_vals) {
        ifstream file(path);
        if (!file.is_open()) {
            cerr << "Khong mo duoc file " << path << "\n";
            return false;
        }

        string line;
        while (getline(file, line)) {
            // Xu ly gia tri missing / o trong
            if (line.empty() || line == "NA" || line == "NaN" || line == "nan") {
                continue; 
            }
            x_vals.push_back(stod(line));
        }

        file.close();
        return true;
    }

    // MODULE 3: Tinh P(x) bang thuat toan Horner
    double evalHorner(vector<double> a, double x) {
        if (a.empty()) return 0.0;
        
        int n = a.size() - 1;
        double res = a[n];
        for (int i = n - 1; i >= 0; i--) {
            res = res * x + a[i];
        }
        return res;
    }

    // MODULE 4: Ghi ket qua ra output.csv
    bool writeOutput(string path, const vector<double>& x_vals, const vector<double>& a) {
        ofstream file(path);
        if (!file.is_open()) {
            cerr << "Khong ghi duoc file " << path << "\n";
            return false;
        }

        // Ghi dong header
        file << "x,P(x)\n";

        // Ghi du lieu
        file << fixed << setprecision(6);
        for (double x : x_vals) {
            double px = evalHorner(a, x);
            file << x << "," << px << "\n";
        }

        file.close();
        return true;
    }

    // MAIN
    int main() {
        int degree = 0;
        vector<double> a;
        vector<double> x_vals;

        if (!readPoly("poly.csv", degree, a)) {
            return 1;
        }

        if (!readXValues("x_values.csv", x_vals)) {
            return 1;
        }

        if (!writeOutput("output.csv", x_vals, a)) {
            return 1;
        }

        return 0;
    }
    \end{lstlisting}

\subsection{Kế hoạch kiểm thử (Test Cases)}

    \textbf{Test Case 1 (Thông thường - Bậc 2)}
    
    $P(x) = 1 + 2x + 3x^2$ (degree=2)
    
    \textbf{Input (poly.csv):}
    \begin{lstlisting}
    2,1.0,2.0,3.0\end{lstlisting}
    \textbf{Input (x\_values.csv):}
    \begin{lstlisting}
    1.5
    -2.0
    NaN\end{lstlisting}
    \textbf{Expected Output (output.csv):}
    \begin{lstlisting}
    x,P(x)
    1.500000,10.750000
    -2.000000,9.000000\end{lstlisting}

    \vspace{0.5cm}

    \textbf{Test Case 2 (Case biên -- Hàm hằng, degree = 0)}
    
    $P(x) = 5.5$
    
    \textbf{Input (poly.csv):}
    \begin{lstlisting}
    0,5.5\end{lstlisting}
    \textbf{Input (x\_values.csv):}
    \begin{lstlisting}
    10.0
    0.0\end{lstlisting}
    \textbf{Expected Output (output.csv):}
    \begin{lstlisting}
    x,P(x)
    10.000000,5.500000
    0.000000,5.500000\end{lstlisting}

    \vspace{0.5cm}
    
    \textbf{Test Case 3 (Case biên -- x\_values rỗng)}

    \textbf{Input (poly.csv):}
    \begin{lstlisting}
    3,1,0,0,5\end{lstlisting}
    \textbf{Input (x\_values.csv):}
    \begin{lstlisting}
    \end{lstlisting}
    (File trống)
    
    \textbf{Expected Output (output.csv):}
    \begin{lstlisting}
    x,P(x)\end{lstlisting}